\documentclass{umemoria}

\depto{Departamento de Ingeniería Civil Industrial}
\author{Mauricio Valenzuela Corvalán}
\title{Fragilidad Bancaria Sistémica y Riesgo de Cola del Crecimiento: Determinantes del Spread Soberano en Economías Emergentes}

\tesis{Magíster en Economía Aplicada}
\memoria{Ingeniero Civil Industrial}

\guia{Juan Francisco Martínez Sepúlveda}
\coguia{Ronald Fischer}
\comision{Patricio Valenzuela, Alejandro Corvalán}
% TODO (completar antes de la entrega final — no se puede inventar):
% \auspicio{Fuente de financiamiento, si corresponde}
% \anho{2026}  % año del examen de grado; por defecto usa el año actual

% Paquetes adicionales requeridos por los capítulos (no incluidos por umemoria.cls)
\usepackage{booktabs}
\usepackage{natbib}
\usepackage{setspace}
\usepackage{mdframed}
\usepackage{longtable}
\usepackage{array}
\usepackage{pdflscape}   % las baterías de regresiones van apaisadas
\usepackage{tikz}        % diagrama del mecanismo conceptual (Sección 1.3)
\usetikzlibrary{arrows.meta,positioning}

% Rutas de figuras
\graphicspath{{imagenes/}}

% Estadístico t bajo el coeficiente en las tablas de regresión (Cap. 2)
\newcommand{\tst}[1]{\ensuremath{\scriptstyle(#1)}}

% Macros de notación compartidas por ambos capítulos y el anexo matemático
\newcommand{\JLoss}{\mathrm{JLoss}}
\newcommand{\Spread}{\mathrm{spread}}
\newcommand{\GaR}{\mathrm{GaR}}
\newcommand{\PD}{\mathrm{PD}}
\newcommand{\ES}{\mathrm{ES}}
\newcommand{\HHI}{\mathrm{HHI}}
\newcommand{\Var}{\mathrm{Var}}
\newcommand{\dd}{\,\mathrm{d}}

% Absorbe micro-desbordes de caja en prosa y bibliografía (URLs, compuestos largos)
\setlength{\emergencystretch}{3em}

% Cada capítulo lleva sus referencias en línea con \begin{thebibliography} y su
% propio encabezado \section*{Referencias del capítulo}. Se suprime el
% \chapter*{Bibliografía} automático (y su salto de página) que la clase book
% inserta al abrir el entorno, para no duplicar el título.
\makeatletter
\renewenvironment{thebibliography}[1]
  {\list{\@biblabel{\@arabic\c@enumiv}}%
        {\settowidth\labelwidth{\@biblabel{#1}}%
         \leftmargin\labelwidth \advance\leftmargin\labelsep
         \@openbib@code \usecounter{enumiv}\let\p@enumiv\@empty
         \renewcommand\theenumiv{\@arabic\c@enumiv}}%
   \sloppy\clubpenalty4000 \@clubpenalty\clubpenalty
   \widowpenalty4000 \sfcode`\.\@m}
  {\def\@noitemerr{\@latex@warning{Empty `thebibliography' environment}}\endlist}
\makeatother

% Cadenas PDF (marcadores / índices con hipervínculos): evita los avisos de hyperref
% "Token not allowed in a PDF string" por la notación matemática en títulos y \caption.
\pdfstringdefDisableCommands{%
  \def\({}\def\){}\def\text#1{#1}\def\mathrm#1{#1}\def\textit#1{#1}\def\textbf#1{#1}%
  \def\hat#1{#1}\def\boldsymbol#1{#1}\def\mathbf#1{#1}%
  \def\JLoss{JLoss}\def\Spread{spread}\def\GaR{GaR}\def\PD{PD}\def\ES{ES}\def\HHI{HHI}\def\Var{Var}%
  \def\times{ x }\def\Rightarrow{=>}\def\to{->}\def\equiv{=}\def\leq{<=}\def\geq{>=}%
  \def\alpha{alpha}\def\beta{beta}\def\theta{theta}\def\rho{rho}\def\gamma{gamma}\def\delta{delta}%
  \def\sqrt#1{raiz(#1)}\def\frac#1#2{#1/#2}\def\partial{d}%
}

\begin{document}

\frontmatter

\maketitle

\begin{resumen}
Esta tesis investiga, sobre un único panel de trece economías emergentes construido con un motor de datos propio, por qué la coincidencia de fragilidad bancaria sistémica y riesgo a la baja del crecimiento amplifica el spread soberano por encima de la suma de sus efectos individuales. La investigación retoma y precisa la hipótesis con la que se aprobó su tema en diciembre de 2025 (Informe de Taller de Tesis I): que el mercado penaliza de forma no lineal la coexistencia de fragilidad bancaria y vulnerabilidad macroeconómica.

Esta tesis estima, sobre \textbf{un único panel} de trece economías emergentes con la fragilidad bancaria construida desde la terminal Bloomberg (muestra de estimación 2004Q1--2026Q1, $N=721$; $N=614$ bajo la especificación con controles), el término de interacción entre la fragilidad bancaria sistémica —medida mediante la métrica \textit{Joint Loss} de Chari et al. (2024), construida a partir de probabilidades de incumplimiento de tipo Merton agregadas por punto de silla, con balances bancarios y precios accionarios de Bloomberg— y el riesgo a la baja del crecimiento —medido mediante \textit{Growth-at-Risk}, estimado por regresión cuantílica de panel siguiendo el marco de \textit{vulnerable growth} de Adrian, Boyarchenko y Giannone (2019)—, sobre el spread soberano medido, siguiendo a Chari et al.\ (2024), por el índice EMBI Global Diversified de J.P.\ Morgan (con el CDS a cinco años como robustez). El esqueleto de la evidencia es una batería de veinticuatro regresiones —cuatro modelos anidados ($JLoss$; $D=-GaR$; ambos; ambos más la interacción) por tres estructuras de efectos fijos por dos muestras—, estructurada al modo de Chari et al.\ (2024). Los canales de nivel están bien identificados en las doce especificaciones: la fragilidad bancaria ($\hat\beta_1>0$, siempre significativa al 1\%) y el riesgo de cola del crecimiento (coef.\ de $D=-GaR$ $>0$, equivalentemente $\hat\beta_2<0$) elevan cada uno el spread soberano. El coeficiente de interacción $\beta_3$ ($JLoss\times D$, $D=-GaR$; positivo bajo la hipótesis de amplificación, $\beta_3=-\theta$) tiene el signo predicho por el multiplicador de \textit{doom loop} pero \textbf{no es significativo sobre la muestra completa} ($\hat\beta_3=+0{,}16$, $p=0{,}26$; invariante a los efectos fijos y a la métrica del spread —EMBI y CDS arrojan el mismo coeficiente sobre la misma submuestra). Esa no significancia es un promedio de regímenes. La interacción \textbf{sí es positiva y significativa} en dos dimensiones: \textbf{transversal} —$\hat\beta_3=+0{,}47$ ($p=0{,}023$) en el núcleo de once economías emergentes de financiamiento externo, diluyéndose al añadir dos economías de deuda local profunda— y \textbf{temporal} —$\hat\beta_3\approx+1{,}2$ ($p<0{,}01$ bajo las tres estructuras de efectos fijos) al excluir los trimestres de la crisis financiera global y la pandemia, cuando los respaldos oficiales desacoplan el canal doméstico. Una especificación más exigente, que en vez de excluir esos trimestres interactúa la complementariedad con un vector de crisis manteniendo toda la muestra, confirma y precisa el resultado ($\hat\beta_3\approx+0{,}8$ fuera de crisis) y aporta un test de falsación directo del mecanismo: la complementariedad \textbf{se anula} bajo los episodios con respaldo oficial masivo (crisis financiera global y pandemia; $\hat\beta_3+\hat\beta_4\approx0$) pero \textbf{sobrevive} en el estrés emergente de 2015--2016, que no tuvo ese respaldo ($\hat\beta_3+\hat\beta_4\approx+1$, $p<0{,}001$). Combinando la dimensión transversal con el corte temporal simple, $\hat\beta_3=+0{,}94$ ($p=0{,}004$). Un modelo de umbral de Hansen corrobora la no linealidad, y una reconstrucción del $GaR$ sin el rendimiento soberano que lo alimenta descarta que estos resultados sean mecánicamente circulares con el EMBI. Las dos dimensiones apuntan al mismo mecanismo: el multiplicador de \textit{doom loop} se traslada al \emph{precio} de la deuda soberana cuando el mercado espera que el pasivo contingente bancario recaiga sobre un soberano \textbf{no respaldado} cuya deuda se fija al margen por inversores sensibles al riesgo —lo que no ocurre ni en las economías de deuda local profunda ni en los episodios de crisis con respaldo oficial—. La complementariedad fragilidad--crecimiento sobre el riesgo soberano se revela, por tanto, como un fenómeno \textbf{condicional}, no como una regularidad universal de los mercados emergentes.

En conjunto, esta tesis (i)~articula como un solo mecanismo dos canales que la literatura trataba por separado —el nexo banca-soberano y el riesgo de cola del crecimiento—; (ii)~muestra que su complementariedad —el multiplicador de \textit{doom loop}— aparece en los precios soberanos de forma condicional, con un mecanismo económico que explica por qué se apaga en las economías de deuda local profunda y en los episodios de crisis con respaldo oficial masivo; (iii)~deja dos advertencias sobre la medición (la variable dependiente EMBI o CDS no altera el resultado; la métrica de fragilidad está censurada por arriba y debe leerse de forma ordinal); y (iv)~lo hace sobre un \textit{pipeline} reproducible. El hilo común es la delimitación precisa de qué está establecido y qué no —una contribución a la credibilidad de la evidencia, no solo a su volumen—. El Anexo~\ref{chap:anexoA} tabula la procedencia de todas las series.
\end{resumen}

% Índice hasta subsección: los subsubsections abultan el índice sin aportar navegación
\setcounter{tocdepth}{2}
\tableofcontents
\listoftables
\listoffigures

\mainmatter

\input{investigacion}

\begin{appendices}
\chapter{Procedencia de los datos}
\label{chap:anexoB}

\noindent Este anexo tabula el origen de todas las series usadas en la investigación empírica
(Capítulo~\ref{chap:paper2}) y en la prueba de las predicciones distintivas del modelo
(Capítulo~\ref{chap:paper1}, Sección~\ref{sec:datos_reales}). El panel se construye sobre el
principio de \textbf{homogeneidad de fuente}: todos los insumos de mercado provienen de la
terminal Bloomberg, extraídos con un mismo protocolo para las veinte economías del universo
objetivo; solo los componentes reales del índice de condiciones financieras que alimenta
$GaR$ y los controles macroeconómicos domésticos provienen de estadísticas oficiales, dado
que Bloomberg no es fuente primaria de cuentas nacionales. El código de extracción y
consolidación reside en \texttt{1\_Codigo/Bloomberg\_extraction/} y \texttt{1\_Codigo/Panel/bbg/}.

\vspace{1em}

\begin{footnotesize}
\setlength{\tabcolsep}{3.5pt}
\renewcommand{\arraystretch}{1.15}
\newcommand{\Lc}{\raggedright\arraybackslash\hyphenpenalty=50\relax\hspace{0pt}}
\begin{longtable}{%
  >{\Lc}p{2.15cm}%
  >{\Lc}p{2.7cm}%
  >{\Lc}p{1.8cm}%
  >{\Lc}p{2.85cm}%
  >{\Lc}p{1.75cm}%
  >{\Lc}p{2.5cm}}
\caption{Procedencia, definición y transformación de cada serie del panel.}
\label{tab:datos}\\
\toprule
\textbf{Serie} & \textbf{Definición} & \textbf{Fuente} & \textbf{Identificador} & \textbf{Frec.\ nativa} & \textbf{Transformación} \\
\midrule
\endfirsthead
\multicolumn{6}{l}{\footnotesize\itshape Tabla~\ref{tab:datos} (continuación)}\\
\toprule
\textbf{Serie} & \textbf{Definición} & \textbf{Fuente} & \textbf{Identificador} & \textbf{Frec.\ nativa} & \textbf{Transformación} \\
\midrule
\endhead
\midrule
\multicolumn{6}{r}{\footnotesize\itshape continúa en la página siguiente}\\
\endfoot
\bottomrule
\endlastfoot

\multicolumn{6}{l}{\textbf{\textit{Variable dependiente}}}\\
\addlinespace
EMBI Global Diversified & Diferencial del índice de bonos soberanos de mercados emergentes de J.P.\ Morgan frente a la curva del Tesoro estadounidense (\textbf{variable dependiente principal}, siguiendo a \citealp{Chari2024}) & J.P.\ Morgan (serie \texttt{embi.xlsx}) & spread EMBI Global por país & Diaria & Media trimestral, puntos básicos. Sin dato $\Rightarrow$ la observación no entra a la regresión de referencia \\
CDS soberano 5Y & Prima del \textit{credit default swap} soberano a cinco años, en dólares, senior no garantizado (\textbf{serie de robustez}) & Bloomberg & \texttt{<emisor> CDS USD SR 5Y Corp}; curva Markit \texttt{CCHIN1U5 Curncy} para China & Diaria & Media trimestral, puntos básicos. Correlación 0,86 con el EMBI en la submuestra común \\
\addlinespace
\multicolumn{6}{l}{\textbf{\textit{Fragilidad bancaria --- $JLoss$}} (métrica compuesta; motor \texttt{jloss\_engine.py})}\\
\addlinespace
Activos totales & Activos totales del banco & Bloomberg & \texttt{BS\_TOT\_ASSET} & Trimestral & Por banco cotizado; ajuste de periodicidad \texttt{ACTUAL} \\
Patrimonio & Patrimonio total contable & Bloomberg & \texttt{TOTAL\_EQUITY} & Trimestral & --- \\
Deuda emitida (LP) & \textit{Long-term borrowings} (bonos, letras, deuda subordinada); define el largo plazo del punto de incumplimiento & Bloomberg & \texttt{BS\_LT\_BORROW} & Trimestral & $D^{*}=$ (pasivo total $-$ bonos) $+ \tfrac12$ bonos \\
Ingreso neto, ingresos, caja & Insumos del $z$-score y de controles de coherencia contable & Bloomberg & \texttt{NET\_\allowbreak INCOME}, \texttt{SALES\_\allowbreak REV\_\allowbreak TURN}, \texttt{BS\_\allowbreak CASH\_\allowbreak NEAR\_\allowbreak CASH\_\allowbreak ITEM} & Trimestral & --- \\
Capital. bursátil & Valor de mercado del patrimonio, cotización local (nunca ADR), verificada contra \texttt{EQY\_\allowbreak FUND\_\allowbreak CRNCY} & Bloomberg & \texttt{CUR\_\allowbreak MKT\_\allowbreak CAP}, o \texttt{PX\_\allowbreak LAST} $\times$ \texttt{BS\_\allowbreak SH\_\allowbreak OUT} & Diaria & Último valor del trimestre \\
Retornos accionarios & Retornos logarítmicos diarios del precio local & Bloomberg & \texttt{PX\_LAST} & Diaria & $\sigma_E=\sqrt{\sum_t r_t^2}$ por banco-trimestre (varianza realizada) \\
Índice bursátil país & Factor sistémico único de la agregación de pérdidas & Bloomberg & índice de referencia por país (\texttt{COUNTRY\_INDEX}) & Diaria & Retornos $\to$ correlación $\rho_i$ banco-factor \\
\addlinespace
\multicolumn{6}{p{15cm}}{\footnotesize\textit{Parámetros del motor (idénticos a \citealp{Chari2024}):} LGD $=0{,}45$; $\rho=0{,}4$; 7 nodos de cuadratura Gauss--Hermite; factor sistémico $N(0,\,0{,}16\%)$; percentil 99; cota de pérdidas $[0{,}010,\,0{,}048]$; horizonte 1 trimestre.}\\
\addlinespace
\multicolumn{6}{l}{\textbf{\textit{Riesgo de cola --- $GaR$}} (regresión cuantílica de panel; motor \texttt{gar\_engine.py}, \texttt{fci\_engine.py})}\\
\addlinespace
PIB real & Crecimiento trimestral del producto (variable dependiente de la regresión cuantílica) & Estadística nacional (banco central) & \texttt{GDP\_*} & Trimestral & Tasa de variación \\
IPC & Componente del FCI (bloque de deuda / expectativas) & Estadística nacional & \texttt{CPI\_*} & Mensual & Estand.\ expansiva \\
Índice accionario & Componente del FCI (bloque accionario) & Bolsa nacional & \texttt{STX\_*} & Mensual & Estand.\ expansiva \\
Tipo de cambio real ef. & Componente del FCI (bloque cambiario) & BIS / fuentes nacionales & \texttt{rEER\_*} & Mensual & Estand.\ expansiva \\
Rendimiento soberano 10Y & Componente del FCI (bloque de deuda soberana) & Mercado / estadística nacional & \texttt{Ryr\_*} & Mensual & Estand.\ expansiva \\
VIX (ortogonalizado) & Condiciones financieras globales en la regresión cuantílica & Bloomberg & \texttt{VIX Index} & Diaria & Ortogonalizado respecto de las condiciones locales; trimestral \\
\addlinespace
\multicolumn{6}{l}{\textbf{\textit{Factores financieros globales}}}\\
\addlinespace
VIX & Índice de volatilidad implícita S\&P 500 & Bloomberg & \texttt{VIX Index} & Diaria & Media trimestral \\
Tasa del Tesoro 10Y & Rendimiento del bono del Tesoro estadounidense a 10 años & Bloomberg & \texttt{USGG10YR Index} & Diaria & Media trimestral; en logaritmo \\
\textit{Spread} HY EE.UU. & \textit{Option-adjusted spread} del \textit{high yield} corporativo estadounidense (equivale a FRED \texttt{BAMLH0A0HYM2}) & Bloomberg & \texttt{LF98OAS Index} & Diaria & Media trimestral; en logaritmo \\
\textit{Spread} on/off-run & Diferencial de liquidez del Tesoro (instrumento \textit{shift-share} principal) & Serie previa del proyecto & \texttt{global\_controls\_} \texttt{quarterly.csv} & Trimestral & En logaritmo \\
Dólar efectivo amplio & Tipo de cambio efectivo nominal de EE.\ UU., canasta de 64 economías (instrumento \textit{shift-share} secundario) & BIS, \textit{Effective Exchange Rate Indices} & \texttt{WS\_\allowbreak EER}, \texttt{M.N.B.US} & Mensual & Media trimestral; en logaritmo \\
\addlinespace
\multicolumn{6}{l}{\textbf{\textit{Controles macroeconómicos domésticos}} (reconstruidos para todos los países del panel)}\\
\addlinespace
Deuda gob.\ general/PIB & Deuda bruta del gobierno general, \% del PIB & FMI, \textit{World Economic Outlook} & \texttt{GGXWDG\_NGDP} & Anual & Interpolación lineal a trimestral \\
Balance fiscal/PIB & Endeudamiento neto del gobierno general, \% del PIB & FMI, \textit{World Economic Outlook} & \texttt{GGXCNL\_NGDP} & Anual & Interpolación lineal a trimestral \\
Reservas/PIB & Reservas internacionales totales sobre PIB corriente & Banco Mundial & \texttt{FI.\allowbreak RES.\allowbreak TOTL.\allowbreak CD} $\div$ \texttt{NY.\allowbreak GDP.\allowbreak MKTP.\allowbreak CD} & Anual & Interpolación lineal a trimestral \\
Cuenta corriente/PIB & Saldo de cuenta corriente, \% del PIB & Banco Mundial & \texttt{BN.\allowbreak CAB.\allowbreak XOKA.\allowbreak GD.\allowbreak ZS} & Anual & Interpolación lineal a trimestral \\
Inflación interanual & Variación del IPC a 12 meses & Serie de IPC del FCI; respaldo Banco Mundial & \texttt{FP.\allowbreak CPI.\allowbreak TOTL.\allowbreak ZG} & Mensual & Variación a 4 trimestres \\
Tipo de cambio real ef. & Índice de tipo de cambio real efectivo (2010 = 100) & BIS; series locales & \texttt{rEER\_*} & Mensual & Último valor del trimestre \\
\addlinespace
\multicolumn{6}{l}{\textbf{\textit{Concentración bancaria e instituciones}}}\\
\addlinespace
$HHI$ (concentración 3 bancos) & Activos de los tres mayores bancos comerciales sobre activos del sistema, \% & Banco Mundial, \textit{Global Financial Development Database} & \texttt{GFDD.\allowbreak OI.01} & Anual & Nivel estructural $=$ mediana por país; serie anual como robustez \\
$HHI_q$ (concentración trimestral) & Activos de los tres mayores bancos \emph{cotizados} sobre activos cotizados totales, \% (mismos bancos que $\JLoss$) & Bloomberg (balances, \texttt{tot\_\allowbreak asset}) & motor \texttt{p6\_\allowbreak concentracion\_\allowbreak trimestral.py} & Trimestral & CR3 y HHI directos; NaN si $<$3 bancos cotizados; corr.\ $\approx$0,5 con GFDD \\
Rating soberano S\&P & Calificación de largo plazo en moneda extranjera & Bloomberg & \texttt{RTG\_\allowbreak SP\_\allowbreak LT\_\allowbreak FC\_\allowbreak ISSUER\_\allowbreak CREDIT} & --- & \textit{Snapshot} (no serie); escala 21--1 \\
Calidad de gobierno (WGI) & Índice de \textit{Worldwide Governance Indicators} (triple interacción institucional) & Banco Mundial & \texttt{instituciones.\allowbreak csv} (provisional) & Anual & Estandarizado; marcado provisional \\

\end{longtable}
\end{footnotesize}

\vspace{1em}

\noindent \textbf{Economías del panel y su cobertura.} El panel incluye toda economía
con $JLoss$ disponible. La muestra de estimación del mecanismo (observaciones con EMBI $+$
$JLoss$ $+$ $GaR$ simultáneamente) reúne \textbf{trece países}: nueve con EMBI de historia
continua (Brasil, Chile, China, Colombia, Malasia, México, Perú, Polonia, Turquía), India
con 54 trimestres (EMBI desde 2012Q4), y Filipinas, Indonesia y Sudáfrica con once cada uno
(EMBI desde 2023Q3, aunque su CDS es de historia larga). Argentina, Egipto y Rusia tienen
$JLoss$ pero no $GaR$; Pakistán tiene $JLoss$ y $GaR$ pero no EMBI. Corea del Sur, Bulgaria
y Hungría se excluyen porque su $JLoss$ no es una medida válida a nivel de país —número de
bancos cotizados insuficiente— (Capítulo~\ref{chap:paper2}, Sección~5.1). La
Figura~\ref{fig:cobertura} traza esta cobertura.

\vspace{1ex}

\noindent La Tabla~\ref{tab:diagpais} resume, por país de la muestra de estimación, el número
mediano de bancos cotizados que alimentan $JLoss$, los trimestres marcados con menos bancos
que el mínimo requerido, y la mediana y el máximo de la serie de $JLoss$. Deja ver dos rasgos
que el Capítulo~\ref{chap:paper2} discute: (i) nueve de los trece países tienen un $JLoss$
mediano entre 2,2 y 2,6 —la variación transversal de la fragilidad la aportan Turquía, India,
China y Brasil—; y (ii) Colombia y México descansan sobre pocas instituciones en una fracción
de los trimestres, aunque sin llegar al caso de Bulgaria (un banco) o Hungría (dos, todos los
trimestres bajo el mínimo) que motivaron su exclusión del panel.

\begin{table}[htbp]
\centering
\footnotesize
\caption{Diagnóstico de $JLoss$ por país (muestra de estimación).}
\label{tab:diagpais}
\begin{tabular}{lrrrr}
\toprule
País & Bancos (mediana) & Trim.\ bajo mínimo & $JLoss$ mediana & $JLoss$ máx. \\
\midrule
Brasil       & 8  & 0  & 4,0  & 15,5 \\
Chile        & 4  & 0  & 2,2  & 11,6 \\
China        & 9  & 2  & 3,8  & 29,0 \\
Colombia     & 3  & 23 & 2,4  & 16,7 \\
Filipinas    & 5  & 0  & 2,9  & 4,3  \\
India        & 11 & 6  & 7,5  & 23,2 \\
Indonesia    & 4  & 0  & 2,3  & 3,7  \\
Malasia      & 6  & 0  & 2,2  & 5,9  \\
México       & 4  & 13 & 2,4  & 13,9 \\
Perú         & 3  & 6  & 2,2  & 8,7  \\
Polonia      & 8  & 0  & 2,6  & 15,0 \\
Sudáfrica    & 6  & 0  & 2,4  & 3,9  \\
Turquía      & 8  & 0  & 6,0  & 28,0 \\
\bottomrule
\end{tabular}

\medskip
\footnotesize\raggedright Fuente: \texttt{1\_Codigo/Panel/bbg/diag\_por\_pais\_bbg.csv}
(motor \texttt{jloss\_engine.py}). ``Trim.\ bajo mínimo'' cuenta los trimestres con menos
bancos cotizados que el umbral del motor. Corea del Sur, Bulgaria y Hungría no aparecen:
quedan fuera del panel por $JLoss$ no válido a nivel de país.
\end{table}

\end{appendices}

\end{document}
